% Root and DE independent proof/code audits passed on 2026-11-15.
\begin{statement}{C.35}{Proved}{The order-thirteen character vanishes on the full pentagon group}
The character $\epsilon_{13}$ constructed in C.33 can be
chosen exactly as there so that
\[
                    \epsilon_{13}|_{\widehat{\mathrm{GT}}}=0.
\]
Consequently the quotient of the full gentle group by the
closed normal subgroup generated by the full pentagon group
has a continuous quotient
\[
                    C_{13}\times\prod_{j\geq1}C_2.
\]
The order-thirteen character remains nonzero on global
multiplier one and does not factor through the metabelian
image.  This assertion gives no nonzero arithmetic-trivial
character of the pentagon group itself.
\end{statement}
\begin{proof}
Use the free-fibre representation of the infinitesimal
four-strand braid Lie algebra and the signed pentagon
defect $\mathcal D(q)$ from P.2.  In particular its
vanishing is necessary for the first nonzero term of
an actual pro-$13$ pentagon solution.
For $q\in S_{12}$, let $R_i(q)$ be the coefficient in the
following table:
\[
\begin{array}{c|c|l}
 i&\text{generator to which }\mathcal D(q)\text{ is applied}
       &\text{output word}\\ \hline
1&a&\mathtt{aaaaaaabcbcbc}\\
2&a&\mathtt{acbaacaaabcba}\\
3&b&\mathtt{aaaabacaacabb}\\
4&b&\mathtt{baaacacbaabac}\\
5&a&\mathtt{abbacbaaacaca}
\end{array}                                                     \tag{C35.1}
\]
In the same integral seven-element basis as C.33, exact
integer expansion gives
\[
\begin{split}
698922\,\ell={}&-192329R_1+472R_2-26842R_3\\
               &-107926R_4-11870R_5 ,
\end{split}                                                     \tag{C35.2}
\]
where $\ell=(-4,1,8,-4,-17,14,84)$.
Thus every degree-twelve linear pentagon solution is
annihilated by $\ell$.  The coefficient $698922$ is a
unit modulo $13$, although the argument below only
needs the identity over $\mathbb Q_{13}$.

Here are the exact verification details for (C35.2).
The independent program \texttt{pbc\_weight12\_pentagon.py}
computes a specified coefficient by contracting adjacent
output-letter pairs.  If a linear derivation $D$ takes
a letter $v$ to a two-letter polynomial, then
\[
 [w]D(P)=
 \sum_{j,v}[w_jw_{j+1}]D(v)\,
       [w_1\cdots w_{j-1}v w_{j+2}\cdots w_m]P.                 \tag{C35.3}
\]
Iterating this exact identity computes the coefficient
of each associative operator word in $q(D,E)$.
The program first compares nineteen independently
computed coefficients with the full degree-seven
pentagon expansion retained in P.2, including its
nonzero coefficient $-16$.  It then finds the five
rows modulo $65521$, recomputes each selected row
over $\mathbb Z$, and verifies (C35.2) on every one
of the seven integral basis vectors.
The entire five-by-seven integer matrix, its coefficient
labels, source hash and exact row combination are retained
in the JSON certificate.  The modular search alone is
not used as a proof of the identity.

It remains to compare this pure-degree certificate with
an arbitrary actual pentagon element, whose first nonzero
term need not have degree twelve.
Let $G,U,\mathfrak u,t$ be precisely the actual local
quotient, its multiplier-one subgroup, its Lazard lattice
and its tame Galois element from C.33.
Let $P$ be the image of the full pentagon group in $G$,
and put $P_1=P\cap U$.
It is a closed subgroup of the finitely generated
nilpotent pro-$13$ group $U$ and is finitely generated;
for example intersect its finite central filtration
with the finite-rank $\mathbb Z_{13}$ graded pieces.
The subgroup part of the Lazard correspondence identifies
$\mathfrak p=\log P_1$ with an actual Lie sublattice of
$\mathfrak u$.

Crucially, $t$ belongs to the Galois image and hence to
$P$.  It normalizes $P_1$ and preserves $\mathfrak p$.
Therefore the integral projector
\[
                     e_{12}=\frac1{12}
                         \sum_{j=0}^{11}\operatorname{Ad}(t)^j
                                                               \tag{C35.4}
\]
preserves $\mathfrak p$ as well.
For any $h\in P_1$, the element
$v=e_{12}\log h$ belongs to
$\mathfrak p\cap\mathfrak u_{12}$.
Its exponential is a genuine element of $P_1$,
and hence is the image of an actual full pentagon pair.
All its local conjugator terms of degrees below twelve
are zero.  Its degree-twelve term $q$ therefore satisfies
the linear pentagon, with no contribution from lower
terms or their higher corrections.
By (C35.2), $\ell(q)=0$.
The character $\theta$ on $U$ in C.33 is precisely the
normalized reduction of $\ell(e_{12}\log h)$.
We have proved
\[
                             \theta(P_1)=0.                    \tag{C35.5}
\]

Let $\widetilde\theta:G\to\mathbb F_{13}$ be the
continuous extension used in C.33.  Equation (C35.5)
shows that its restriction to $P$ factors through
the local multiplier.
For an even more direct comparison, given $g\in P$
choose a Galois element $\sigma\in P$ with
$\lambda_{13}(\sigma)=\lambda_{13}(g)$.
Then $g\sigma^{-1}\in P_1$, so
$\widetilde\theta(g)=\widetilde\theta(\sigma)$.
The cyclotomic correction in (C33.8) was chosen to
remove this value on every Galois element.
It consequently removes the value on every $g\in P$.
This is the same character $\epsilon_{13}$,
with the same multiplier-one witness.

Finally combine $\epsilon_{13}$ with the independent
quadratic characters of C.30.  Each finite joint
projection is onto by coprime orders and the stated
quadratic independence.  Compactness gives the displayed
product quotient.  The character is zero on the entire
pentagon group, so it makes no assertion about that
group's own arithmetic-trivial character space.
\end{proof}
